\section{Dimension-free matrix inequalities}\label{sec:matrix}

The finite-dimensional part of the argument has two tasks: to bound a
cyclic product by positive quadratic costs, and to keep those costs
under control when matrices change or are averaged.  We give all the
needed estimates here, including the support and complex-phase
arguments required for arbitrary complex input matrices.
The averaging inequality in Lemma~\ref{lem:joint-convexity} will permit
a Gaussian window to be replaced by a probability average of narrower
windows without paying for their relative widths.

All spaces in this section are finite-dimensional complex Hilbert
spaces.  Matrix inequalities mean inequalities of Hermitian quadratic
forms.  Set
\[
 \Phi(P)=\tr(P^{3/2}),\qquad \supp P=(\ker P)^\perp
 \quad(P\ge0).
\]
The application uses three spaces $H_v$, indexed by
$v\in\mathbb Z/3\mathbb Z$, and maps $W_v:H_{v+1}\to H_v$.
The two matrices incident to vertex $v$ form its \emph{row star}
\begin{equation}\label{eq:row-star}
 R_v=[\,W_v\quad W_{v-1}^*\,]:
 H_{v+1}\oplus H_{v-1}\longrightarrow H_v.
\end{equation}
We use both Gram matrices
\begin{equation}\label{eq:star-grams}
 T_v=R_v^*R_v,\qquad S_v=R_vR_v^*.
\end{equation}
Their nonzero eigenvalues agree, including multiplicities, so
$\Phi(T_v)=\Phi(S_v)$. Our eventual objective in this section is to
bound the cyclic trace formed from the $W_v$ by second-derivative
costs of these three energies. We first develop those costs for an
arbitrary positive matrix $P$.

A matrix is \emph{supported on $P$} if it vanishes on $\ker P$ on
both sides.  The Hilbert--Schmidt inner product is
$\langle K,L\rangle_{\HS}=\tr(K^*L)$.  Write
$\mathsf L_A(X)=AX$ and $\mathsf R_A(X)=XA$.  For Hermitian matrices
$K,L$ supported on $P$, define
\begin{equation}\label{eq:hessian-form}
 \mathcal B_P(K,L)=\frac32\operatorname{Re}
 \left\langle K,
 (\mathsf L_{\sqrt P}+\mathsf R_{\sqrt P})^{-1}L
 \right\rangle_{\HS},
 \qquad b_P(K)=\mathcal B_P(K,K).
\end{equation}
The inverse is taken on matrices acting on $\supp P$, where the
operator is positive definite.  If $P=0$, its only supported direction
is zero and these expressions are zero.

\subsection{The energy and its quadratic costs}

\begin{lemma}[Supported Hessian]\label{lem:hessian}
On the cone of matrices positive definite on a fixed support,
$\Phi$ is smooth and
\[
 d\Phi_P(K)=\frac32\tr(\sqrt P K),\qquad
 d^2\Phi_P(K,L)=\mathcal B_P(K,L).
\]
In an orthonormal eigenbasis for the positive eigenvalues $p_i$ of
$P$,
\begin{equation}\label{eq:hessian-metric}
 b_P(K)=\frac32\sum_{i,j}
 \frac{|K_{ij}|^2}{\sqrt{p_i}+\sqrt{p_j}}.
\end{equation}
In particular, $\mathcal B_P$ is a positive definite symmetric
bilinear form on the real space of supported Hermitian directions,
and
\[
 |\mathcal B_P(K,L)|\le b_P(K)^{1/2}b_P(L)^{1/2}.
\]
\end{lemma}

\begin{proof}
Restrict to the support and put $A=\sqrt P$.  The derivative of the
squaring map at $A$ is $X\mapsto AX+XA$; in an eigenbasis, it
multiplies entry $(i,j)$ by $\sqrt{p_i}+\sqrt{p_j}>0$.
The inverse function theorem gives a smooth positive square root,
whose derivative $X_K$ solves
\[
 AX_K+X_KA=K.
\]
Multiplying by $A$ and taking traces gives
$\tr(PX_K)=\tfrac12\tr(AK)$.  Hence differentiation of
$\Phi(P)=\tr(P\sqrt P)$ yields
$d\Phi_P(K)=\tfrac32\tr(AK)$.  Differentiating once more proves
the Hessian formula.  Formula~\eqref{eq:hessian-metric} follows by
diagonalizing $A$, and proves the remaining assertions.
This is the square-root case of the classical divided-difference
calculus~\cite[Section~1]{Daletskii1957}.
\end{proof}

Both $\Phi$ and $\mathcal B$ are invariant under simultaneous unitary
conjugation of their arguments.  They are also invariant under
extension of every argument by zero: their positive eigenvalues and
supported entries do not change.  These facts will let us compare
Grams on different subspaces.

\begin{lemma}[Order comparison]\label{lem:order}
If $P,Q\ge0$ on the same space, $a\ge1$, and $P\le aQ$, then every
Hermitian $K$ supported on $P$ is supported on $Q$, and
\begin{equation}\label{eq:matrix-order}
 b_P(K)\ge a^{-1/2}b_Q(K).
\end{equation}
\end{lemma}

\begin{proof}
For $x\in\ker Q$ we have $0\le\langle Px,x\rangle
\le a\langle Qx,x\rangle=0$, so $\ker Q\subseteq\ker P$.
This proves the support assertion.

We recall the order facts used in the proof.  Inversion reverses
order for positive definite matrices, as follows by conjugating
$0<X\le Y$ by $X^{-1/2}$, inverting, and conjugating back.  Also
\[
 \sqrt X=\frac1\pi\int_0^\infty
 X(X+tI)^{-1}\,\frac{\dd t}{\sqrt t}\qquad(X\ge0).
\]
Spectral calculus reduces this norm-convergent formula to the scalar
integral.  Since $X(X+tI)^{-1}=I-t(X+tI)^{-1}$, inversion order
shows that square root preserves order. This is the classical
Loewner--Heinz square-root monotonicity; see, for example,
Kwong~\cite[Theorem~1 and Section~2]{Kwong1975}. The supported
Hessian comparison below is derived from these order facts here.

For $\eta>0$,
\[
 P+\eta I\le a(Q+\eta I),\qquad
 \sqrt{P+\eta I}\le\sqrt a\sqrt{Q+\eta I}.
\]
If $A\le B$, then
$\mathsf L_A+\mathsf R_A\le\mathsf L_B+\mathsf R_B$ on
Hilbert--Schmidt space: the difference has quadratic form
$\tr(X^*(B-A)X)+\tr(X^*X(B-A))\ge0$.  Invert these positive
operators to obtain
\[
 b_{P+\eta I}(K)\ge a^{-1/2}b_{Q+\eta I}(K).
\]
Let $\eta\downarrow0$.  In an eigenbasis, every entry of $K$ meeting
the kernel of the corresponding limiting matrix is zero.  Thus
Formula~\eqref{eq:hessian-metric} gives precisely the supported costs
on both sides of \eqref{eq:matrix-order}.
\end{proof}

\begin{lemma}[Inserted Gram bound]\label{lem:insertion}
For a linear map $R:E\to F$ and a Hermitian map $L:F\to F$, the
matrix $R^*LR$ is supported on $R^*R$, and
\begin{equation}\label{eq:insertion-bound}
 b_{R^*R}(R^*LR)
 \le\frac34\lVert L\rVert_{\op}^2\Phi(R^*R).
\end{equation}
\end{lemma}

\begin{proof}
The insertion is Hermitian and annihilates $\ker R=\ker(R^*R)$.
If $R=0$, there is nothing to prove.  Otherwise let $\rho_i>0$ be
the positive singular values of $R$, and let $C$ be the compression
of $L$ to the corresponding left singular vectors.  Then
$\lVert C\rVert_{\op}\le\lVert L\rVert_{\op}$ and
\[
 b_{R^*R}(R^*LR)
 =\frac32\sum_{i,j}
 \frac{\rho_i^2\rho_j^2}{\rho_i+\rho_j}|C_{ij}|^2.
\]
For $x,y>0$,
$x^2y^2/(x+y)\le (xy)^{3/2}/2\le(x^3+y^3)/4$.
Every row and column of $C$ has squared Euclidean norm at most
$\lVert L\rVert_{\op}^2$.  Therefore
\[
 b_{R^*R}(R^*LR)
 \le\frac38\sum_{i,j}(\rho_i^3+\rho_j^3)|C_{ij}|^2
 \le\frac34\lVert L\rVert_{\op}^2\sum_i\rho_i^3,
\]
as required.
\end{proof}

We next show that averaging a Gram and its inserted direction can
only decrease the corresponding cost.  No uniform lower bound on
the positive eigenvalues is required.

\begin{lemma}[Joint convexity under probability averages]
\label{lem:joint-convexity}
Let $(\Omega,\nu)$ be a probability space and $E$ a fixed subspace
of a finite-dimensional complex Hilbert space.  Suppose
$T_\omega\ge0$ and $K_\omega=K_\omega^*$ are measurable matrices,
$\supp T_\omega=E$ and $K_\omega$ is supported on $T_\omega$ almost
everywhere, and
\[
 \int_\Omega\bigl(\lVert T_\omega\rVert_{\op}
                  +\lVert K_\omega\rVert_{\op}\bigr)\,\dd\nu(\omega)
 <\infty.
\]
Then $\overline T=\int T_\omega\,\dd\nu$ has support $E$,
$\overline K=\int K_\omega\,\dd\nu$ is supported on $\overline T$, and
\begin{equation}\label{eq:joint-convexity}
 b_{\overline T}(\overline K)
 \le\int_\Omega b_{T_\omega}(K_\omega)\,\dd\nu(\omega).
\end{equation}
The right side may be infinite.
\end{lemma}

\begin{proof}
If $E=\{0\}$, all matrices and costs vanish.  Otherwise restrict
every matrix to $E$.  For every nonzero $x\in E$, the positive
quantity $\langle T_\omega x,x\rangle$ has positive integral, so
$\overline T$ is positive definite on $E$.  All averages vanish on
$E^\perp$, proving the support assertions.

Put $A_\omega=\sqrt{T_\omega}$ and $A=\int A_\omega\,\dd\nu$.
The integral exists by the stated integrability and Cauchy--Schwarz.
For every vector $x$, the vector-valued Cauchy--Schwarz inequality
gives
\[
 \lVert Ax\rVert^2
 \le\int\lVert A_\omega x\rVert^2\,\dd\nu
 =\langle\overline T x,x\rangle.
\]
Consequently $A^2\le\overline T$, and square-root order,
proved in Lemma~\ref{lem:order}, implies
\begin{equation}\label{eq:square-root-average}
 \sqrt{\overline T}\ge\int\sqrt{T_\omega}\,\dd\nu.
\end{equation}
For a positive definite selfadjoint operator $\mathsf M$ on a real
Hilbert space, completing the square proves
\begin{equation}\label{eq:inverse-variational}
 \langle K,\mathsf M^{-1}K\rangle
 =\sup_X\{2\langle K,X\rangle-\langle X,\mathsf M X\rangle\}.
\end{equation}
Apply this identity on Hermitian matrices with real Hilbert--Schmidt
inner product.  The operators
$\mathsf L_{\sqrt T}+\mathsf R_{\sqrt T}$
preserve that real space.  By \eqref{eq:square-root-average}, for
each Hermitian $X$ the expression inside the supremum for
$(\overline T,\overline K)$ is at most the average of the
corresponding expressions for $(T_\omega,K_\omega)$.
Taking the supremum and then bounding the supremum of an integral
by the integral of the pointwise suprema gives
\[
 b_{\overline T}(\overline K)
 \le\int b_{T_\omega}(K_\omega)\,\dd\nu.
\]
Every inverse in this argument is taken on matrices acting on $E$;
uniform invertibility over $\omega$ was never assumed.
\end{proof}

\subsection{Changes that alter the rank}

The Hessian was defined on a fixed support, but switching matrix
entries on or off can change that support.  The following
first-derivative formula remains valid across such changes.  It is
a special case of the gradient formula for unitarily invariant
matrix functions~\cite[Theorem~3.1]{Lewis1995}.

\begin{lemma}[Rank changes and entrywise derivatives]
\label{lem:rank-change}
For rectangular complex matrices $R_-,R_+$ of the same size, set
$R(z)=(1-z)R_-+zR_+$ and $D=R_+-R_-$.  Then
\begin{equation}\label{eq:rank-change-derivative}
 \Phi(R_+^*R_+)-\Phi(R_-^*R_-)
 =3\int_0^1\operatorname{Re}
 \langle R(z)\sqrt{R(z)^*R(z)},D\rangle_{\HS}\,\dd z.
\end{equation}
For every rectangular matrix $R$ and every entry $(i,j)$,
\begin{equation}\label{eq:entry-gradient}
 |(R\sqrt{R^*R})_{ij}|
 \le\lVert\operatorname{row}_iR\rVert_2
       \lVert\operatorname{col}_jR\rVert_2.
\end{equation}
\end{lemma}

\begin{proof}
For $\eta>0$, Lemma~\ref{lem:hessian} gives
\[
 \frac{\dd}{\dd z}\tr((R(z)^*R(z)+\eta I)^{3/2})
 =3\operatorname{Re}
 \langle R(z)\sqrt{R(z)^*R(z)+\eta I},D\rangle_{\HS}.
\]
Spectral calculus implies
$\lVert\sqrt{R^*R+\eta I}-\sqrt{R^*R}\rVert_{\op}\le\sqrt\eta$.
The matrices $R(z)$ stay bounded for $0\le z\le1$, so the
derivatives converge uniformly as $\eta\downarrow0$, and the
endpoint energies converge.  Integrating first and passing to the
limit proves \eqref{eq:rank-change-derivative}.

For \eqref{eq:entry-gradient}, Cauchy--Schwarz in the matrix product
bounds the entry by the norm of row $i$ of $R$ times the norm of
column $j$ of $\sqrt{R^*R}$.  The latter column has squared norm
$(R^*R)_{jj}$, also the squared norm of column $j$ of $R$.
\end{proof}

\subsection{The cyclic product}

Return to the three maps $W_v$ and the row stars
\eqref{eq:row-star}--\eqref{eq:star-grams}. The reason for combining
the two incident matrices is the following estimate: three inserted
star costs control the entire cyclic product.

\begin{lemma}[Mixed trace inequality]\label{lem:mixed-trace}
For Hermitian maps $D_v:H_v\to H_v$, put
$D_* =\max_v\lVert D_v\rVert_{\op}$.  Then
\begin{equation}\label{eq:mixed-trace}
 \left|\tr(D_0W_0D_1W_1D_2W_2)\right|
 \le\frac53D_*\sum_{v=0}^2b_{T_v}(R_v^*D_vR_v).
\end{equation}
The constant is independent of the dimensions and no Gram is
assumed invertible.
\end{lemma}

\begin{proof}
We put the three costs in one quadratic form, bound a cubic trace,
and then recover the complex phase of the cyclic product.

\paragraph{One quadratic form.}
On $H=H_0\oplus H_1\oplus H_2$, define
\[
 M=\begin{pmatrix}
 0&W_0&W_2^*\\
 W_0^*&0&W_1\\
 W_2&W_1^*&0
 \end{pmatrix},\qquad
 D=\operatorname{diag}(D_0,D_1,D_2).
\]
Both are Hermitian and $\lVert D\rVert_{\op}=D_*$.  If $P_v$ is
the orthogonal projection onto $H_v$, put
$A_v=MP_vM$ and $K_v=MP_vDM$.  After permuting blocks and extending
by zero, these are respectively $T_v$ and $R_v^*D_vR_v$.
The matrix $K_v$ is Hermitian because $P_v$ commutes with $D$.
Moreover,
\[
 \ker A_v=\ker(P_vM)\subseteq\ker K_v,
 \qquad 0\le A_v\le M^2.
\]
Thus $K_v$ is supported on $A_v$ and on $M^2$.  Lemma~\ref{lem:order}
and the inequality $\sum_{v=0}^2 b(X_v)\ge\tfrac13 b(\sum_vX_v)$
for a positive quadratic form give
\begin{equation}\label{eq:trace-common-metric}
 \mathcal E:=\sum_v b_{T_v}(R_v^*D_vR_v)
 \ge\sum_v b_{M^2}(K_v)
 \ge\frac13 b_{M^2}(MDM).
\end{equation}

\paragraph{A cubic trace estimate.}
Choose an orthonormal eigenbasis of $M$, with real eigenvalues $m_i$
ordered so that $q_i=|m_i|$ is nonincreasing.  Ties are arbitrary.
In that basis Formula~\eqref{eq:hessian-metric} yields
\begin{equation}\label{eq:trace-Q}
 \mathcal E\ge\frac12Q,\qquad
 Q:=\sum_{i,j:q_i+q_j>0}
 \frac{q_i^2q_j^2}{q_i+q_j}|D_{ij}|^2.
\end{equation}
Any term with one zero $q_i$ is zero; pairs with both zero are
omitted.  Write
\[
 A_i=q_i^3|D_{ii}|^2,\quad
 E_i=q_i\sum_{j>i}q_j^2|D_{ij}|^2,\quad
 A=\sum_iA_i,\quad E=\sum_iE_i.
\]
Hermitian symmetry and $q_j\le q_i$ for $j>i$ imply
\begin{equation}\label{eq:trace-AE}
 Q=\frac A2+
 2\sum_{\substack{i<j\\q_i+q_j>0}}
 \frac{q_i^2q_j^2}{q_i+q_j}|D_{ij}|^2
 \ge\frac A2+E.
\end{equation}
Expand
\[
 \tr((DM)^3)=\sum_{i,j,k}m_im_jm_kD_{ij}D_{jk}D_{ki}
\]
and group triples by the multiplicity of their smallest index $i$.
If it occurs once, its three cyclic positions give
\[
 3m_i\sum_{j,k>i}m_jm_kD_{ij}D_{jk}D_{ki}.
\]
The inner sum is $x^*D_{>i}x$ for
$x=(m_jD_{ji})_{j>i}$, where $D_{>i}$ is the corresponding
compression of $D$.  Since $\lVert D_{>i}\rVert_{\op}\le D_*$,
this contribution has modulus at most $3D_*E_i$.  The argument
allows $j=k$ and retains the signs of the $m_j$.

If the smallest index occurs twice, its contribution is
\[
 3m_i^2D_{ii}\sum_{k>i}m_k|D_{ik}|^2.
\]
By Cauchy--Schwarz and
$\sum_k|D_{ik}|^2=\lVert De_i\rVert^2\le D_*^2$,
\[
 \left|\sum_{k>i}m_k|D_{ik}|^2\right|
 \le D_*\left(\sum_{k>i}q_k^2|D_{ik}|^2\right)^{1/2}.
\]
The contribution is therefore at most $3D_*\sqrt{A_iE_i}$, also
when $q_i=0$.  Finally, if all three indices equal $i$, the modulus
is at most $D_*A_i$.  These cases exhaust all triples.  Summing and
using Cauchy--Schwarz gives
\begin{equation}\label{eq:trace-cubic}
 \begin{aligned}
 |\tr((DM)^3)|
 &\le D_*\bigl(A+3E+3\sqrt{AE}\bigr)\\
 &\le D_*\left(\frac52A+\frac92E\right)
 \le5D_*\left(\frac A2+E\right)
 \le5D_*Q.
 \end{aligned}
\end{equation}

\paragraph{Recovering the complex phase.}
Set $\tau=\tr(D_0W_0D_1W_1D_2W_2)$.  Since $M$ has zero
diagonal blocks, every nonzero closed three-step block product
visits all three vertices.  The three starting vertices in one
orientation give $\tau$, by cyclicity of trace; the reverse
orientation gives $\overline\tau$, because the $D_v$ are Hermitian.
Consequently
\[
 \tr((DM)^3)=6\operatorname{Re}\tau.
\]
When $\tau\ne0$, choose $|z|=1$ so that $z\tau=|\tau|$ and replace
$W_0$ by $zW_0$.  This changes the stars by
\[
 R_0\longmapsto R_0\operatorname{diag}(zI_{H_1},I_{H_2}),\qquad
 R_1\longmapsto R_1\operatorname{diag}(I_{H_2},\overline zI_{H_0}),
\]
and leaves $R_2$ unchanged.  These right multipliers are unitary;
each conjugates the Gram and its insertion by the same unitary.
Thus $\mathcal E$ and $D_*$ are unchanged.  Apply
\eqref{eq:trace-Q} and \eqref{eq:trace-cubic} to the modified
matrices, with their own value of $Q$, to obtain
\[
 6|\tau|\le5D_*Q\le10D_*\mathcal E.
\]
This is \eqref{eq:mixed-trace}.  The case $\tau=0$ is immediate.
\end{proof}

In the applications a diagonal insertion may be complex.  Write
$D_v=D_v^{(0)}+iD_v^{(1)}$ for its real and imaginary diagonal
parts, and put $D_* =\max_v\lVert D_v\rVert_{\op}$.  Expanding
gives eight products with Hermitian diagonals, each bounded by
Lemma~\ref{lem:mixed-trace}.  Both parts have operator norm at most
$D_*$, and each of their six costs occurs four times in the sum.
Thus
\begin{equation}\label{eq:complex-mixed-trace}
 \left|\tr(D_0W_0D_1W_1D_2W_2)\right|
 \le\frac{20}{3}D_*\sum_{v=0}^2\sum_{\epsilon=0}^1
 b_{T_v}(R_v^*D_v^{(\epsilon)}R_v).
\end{equation}
All directions appearing in this bound remain Hermitian.
